\documentclass[a4paper,12pt,parskip=half,headsepline,DIV=12,numbers=noenddot]{scrartcl}

% Loading all packages and configuration present in the preamble file
\usepackage{preamble}

%%%%%%%%%%%%%%%%%%%%%%%%%%%%%%%%%%%%%%%%%%%%%%%%%%%%%%%%%%%%%%%%%%%
%  Special Characters:                                            %
%                                                                 %
%             \& \% \$ \# \_ \{ \}                                %
%             \textasciitilde (~)                                 %
%             \textasciicircum (^)                                %
%             \textbackslash (\)                                  %
%      \glqq Text\grqq{} für Anführungszeichen                    %
%%%%%%%%%%%%%%%%%%%%%%%%%%%%%%%%%%%%%%%%%%%%%%%%%%%%%%%%%%%%%%%%%%%

% ======= BEGIN DES DOKUMENTS =======
\begin{document}

% Definition Header Sections sollen in der Kopfzeile stehen; Kopfzeile mit Unterstrich
\automark[subsection]{section}
\KOMAoptions{headsepline=true}
%\ihead{Kopfzeile innen}
%\chead{Kopfzeile außen}
\ohead{\headmark}

% Definition footer \pagemark steht für Seitennummer
%\ifoot{Fußzeile innen}
%\cfoot{Fußzeile Mitte}
\ofoot{\pagemark}

% % ------------------ Titelseite ------------------
%%%%%%%%%%%%%%%%%%%%%%%%%%%%%%%%%%%%%%%%%%%%%%%%%%%%%%%%%%%%%%%%%%%
%                                                                 %                    
%                    Definition Deckblatt                         %
%                                                                 %
%%%%%%%%%%%%%%%%%%%%%%%%%%%%%%%%%%%%%%%%%%%%%%%%%%%%%%%%%%%%%%%%%%%

% true für Bachelorarbeit / false für Hausarbeit
\newbool{bachelorarbeit}
\setbool{bachelorarbeit}{false}

% Setze Fachbereich
\newcommand{\department}{Fachbereich II \\ Management und Informationssysteme}

% Setze Studiengang
\newcommand{\studyprogram}{Wirtschaftsinformatik B.Sc.}

% Setze Modulname (bachelorarbeit muss false sein)
\newcommand{\modulname}{Tolles Modul}

% Setze Dozent:in (bachelorarbeit muss false sein)
\newcommand{\auditor}{\textbf{Dozent:in:} \> Prof. Dr. Maxi Muster}

% Setze Gutachter:innen (bachelorarbeit muss true sein)
\newcommand{\firstauditor}{\textbf{Erstgutachter:} \> Prof. Dr. Maxi Mustermann}
\newcommand{\secondauditor}{\textbf{Zweitgutachterin:} \> Prof. Dr. Maxi Musterfrau}

% Setze Titel und Untertitel der Abreit 
\newcommand{\thetitle}{Titel}
\newcommand{\thesubtitle}{Untertitel}

% Setze Autor:in und MatNr.
\newcommand{\theauthor}{Robin Mustermensch}
\newcommand{\matriculationnumber}{00000}

% Abstand zwischen Name und MatNr. (siehe Deckblatt)
\newcommand{\myspace}{1.0cm}

% Muss in src/basic_structure/deckblatt.tex einkommentiert werden! 

\newcommand{\secondauthor}{\> Maxi Mustermensch \> MatNr. 00000\\}
\newcommand{\thirdauthor}{\> Maxi Mustermensch \> MatNr. 00000\\}
\newcommand{\fourthauthor}{\> Maxi Mustermensch \> MatNr. 00000\\}
\newcommand{\fifthauthor}{\> Maxi Mustermensch \> MatNr. 00000\\}

% PDF Metadaten
%\hypersetup{pdfinfo={
%Title={\thetitle},
%Author={\theauthor}
%}}

\hypersetup{pdfinfo={
			Title={\thetitle},
			Author={\theauthor}
		}}


%%%%%%%%%%%%%%%%%%%%%%%%%%%%%%%%%
%           Deckblatt           %
%%%%%%%%%%%%%%%%%%%%%%%%%%%%%%%%%
\setmainfont{TeX Gyre Adventor}
\thispagestyle{empty}
\begin{figure}[h!]
	\centering
	\includegraphics[width=0.6\textwidth]{./images/hs-logo.png}
\end{figure}
\begin{center}
	\large{\textbf{\department}}\\
	\large{\textbf{\studyprogram}}\\
	\vspace{1cm}
	\ifbool{bachelorarbeit}{
		\LARGE{\textbf{Bachelorarbeit}}\\
		\large{zur Erlangung des akademischen Grades \\ Bachelor of Science}\\
	}
	{
		\large{\textbf{Modul\\ \modulname}}\\
	}
	\vspace*{\fill}
	\line(1,0){450}\\
	\doublespacing
	\textbf{\Large{\thetitle}}\\
	\textbf{\large{\thesubtitle}}\\
	\line(1,0){450}\\
\end{center}
\vspace*{\fill}
\onehalfspacing
\small{
	\begin{flushleft}
		\begin{tabbing}
			\textbf{Vorgelegt von:} \hspace*{0.8cm}\= \theauthor \hspace*{\myspace}\= MatNr. \matriculationnumber \\

			%%%%%%%%%%%%%%%%%%%%%%%%%%%%%%%%%%%%%%%%%%%%%%
			%                                            %
			%  Hier weitere Autor:innen einkommentieren  %
			%   Müssen hauptdatei.tex definiert sein     %
			%	                                         %
			%%%%%%%%%%%%%%%%%%%%%%%%%%%%%%%%%%%%%%%%%%%%%%		
			%\secondauthor
			%\thirdauthor
			%\fourthauthor
			%\fifthauthor

			\textbf{Vorgelegt am:} \> \today\\
			\ifbool{bachelorarbeit}{
				\firstauditor\\
				\secondauditor\\
			}
			{
				\auditor\\
			}
		\end{tabbing}
	\end{flushleft}}
\setmainfont{TeX Gyre Termes}

\newpage

% Singlespacing (Zeilenabstand) (Default)
\singlespacing
\normalsize

% Abstract falls gewünscht
% \thispagestyle{empty}
% \input{./contents/abstract}
% \newpage

% Inhaltsverzeichnis anzeigen
\pagestyle{empty}
\tableofcontents
\newpage
\pagestyle{headings}

% Header für den Inhalt
\KOMAoptions{headsepline=true}
\ohead{\headmark}

% % ------------------ Inhalt des Dokuments ------------------



\section{Einleitung}
Das hier ist die Einleitung
\section{Erster Teil}
\input{contents/2.erster.tex}
\section{Zweiter Teil}
\input{contents/3.zweiter.tex}
\section{Zusammenfassung}
\input{contents/4.zusammenfassung.tex}
\newpage


% % ------------------ Verzeichnisse ------------------
% Literaturverzeichnis anzeigen
\ohead{Literaturverzeichnis} % Korrektur für Header
\phantomsection
\addcontentsline{toc}{section}{Literaturverzeichnis}
\renewcommand\refname{Literaturverzeichnis}
% \printbibliography
\newpage

% Abbildungsverzeichnis anzeigen
\ohead{\headmark}
\listoffigures
\addcontentsline{toc}{section}{Abbildungsverzeichnis}
\newpage


% Tabellenverzeichnis anzeigen
\listoftables
\addcontentsline{toc}{section}{Tabellenverzeichnis}
\newpage

%Listingverzeichnis anzeigen
\lstlistoflistings
\addcontentsline{toc}{section}{Quellcodeverzeichnis}
\newpage


% Abkürzungsverzeichnis anzeigen
\ohead{Abkürzungsverzeichnis} % Korrektur für Header
\section*{Abkürzungsverzeichnis}
%%%%%%%%%%%%%%%%%%%%%%%%%%%%%%%%%%%%%%%%%%%%%%%%%%%%%%%%%%%%%%%%%%%%%%%%%%%%%%%%%%%%%%%%%
%                             Mit Einrückung im Verzeichnis                             %
%  Alle Abkürzungen müssen mit: \acro{acronym}[shortname]{fullname} eingetragen werden. %
%                         Im Fließtext mit \ac{acronym} aufrufen.                       %
%                 Erster Aufruf: fullname (shortname), danach shortname.                %
%%%%%%%%%%%%%%%%%%%%%%%%%%%%%%%%%%%%%%%%%%%%%%%%%%%%%%%%%%%%%%%%%%%%%%%%%%%%%%%%%%%%%%%%%

\begin{addmargin}[1.5em]{0em}

  \begin{acronym}[HTTPS] %längster shortname muss eingetragen werden.

    \acro{eugh}[EuGH]{Europäischer Gerichtshof}

    \acro{eu}[EU]{Europäische Union}

    \acro{https}[HTTPS]{Hypertext Transfer Protocol Secure}

  \end{acronym}

\end{addmargin}

\addcontentsline{toc}{section}{Abkürzungsverzeichnis}
\newpage

% Kein Header für Anhang (Deckblatt)
\KOMAoptions{headsepline=false}
\ohead{}

% ------------------ Anhang ------------------
%%%%%%%%%%%%%%%%%%%%%%%%%%%%%%%%%
%      Deckblatt Anhang         %
%%%%%%%%%%%%%%%%%%%%%%%%%%%%%%%%%
\phantomsection
\addcontentsline{toc}{section}{Anhang}
\vspace*{\fill}
\begin{center}
  \doublespacing
  \line(1,0){445}\\
  \textbf{\Huge{Anhang}}\\
  \line(1,0){445}
  \singlespacing
  \begin{verbatim}
    \end{verbatim}
\end{center}
\vspace*{\fill}
\newpage

% % Anhang römisch
\renewcommand{\thesection}{\Roman{section}}
\renewcommand{\thesubsection}{\Roman{subsection}}
\setcounter{section}{0}
\counterwithin{table}{subsection}
\counterwithin{figure}{subsection}
% \counterwithin{listing}{subsection}

% % Header Anhang (Inhalt)
\KOMAoptions{headsepline=true}
\ohead{\headmark}
\automark{subsection}

% % Input Anhang
\input{./contents/anhang}
\newpage

%------------------ Selbstständigkeitserklärung ------------------

% Selbstständigkeits Erklärung
\phantomsection
\addcontentsline{toc}{section}{Selbstständigkeitserklärung}

% Header für Erklärung
\ohead{Selbstständigkeitserklärung}

% Input Erklärung
\section*{Selbstständigkeitserklärung}

%\vspace{1cm}
Ich versichere, die von mir vorgelegte Arbeit selbstständig verfasst zu haben. Alle Stellen, die wörtlich oder sinngemäß aus veröffentlichten oder nicht veröffentlichten Arbeiten anderer entnommen sind, 
habe ich als entnommen kenntlich gemacht. Sämtliche Quellen und Hilfsmittel, die ich für die Arbeit benutzt habe, sind angegeben. Die Arbeit habe ich mit gleichem Inhalt bzw. in wesentlichen 
Teilen noch keiner anderen Prüfungsbehörde vorgelegt.

\vspace*{1cm}

\begin{tblr}{X[l] X[l]}
Bremerhaven, den \today & Unterschrift:\\
\end{tblr}

% ======= ENDE DES DOKUMENTS =======
% Leere Abschlussseite
\newpage
\thispagestyle{empty} % erzeugt Seite ohne Kopf- / Fusszeile
\mbox{}

\end{document}